\section{Scope, source snapshot, and result map}
\label{sec:scope}

The source manuscript is Vladimir Reshetnikov's \repo\ project,
under \codefile{Analysis/Polylogarithms/docs/manuscript}. We pinned the
working branch to commit \baseline\ and inspected its chapter sources,
editorial ledger, and the relevant pending research reports
\cite{ProveIt,AngularBaseline,LerchBaseline,CayleyBaseline}.
This matters mathematically: several older reports call $S_4$ a conjecture,
whereas the canonical manuscript now contains an exact proof. Its
specified level-four product-matrix rank conjecture is also already
proved. The $S_6$ depth-two identity is still conjectural, and an earlier
frozen $S_8$ search vector has already been rigorously refuted. None of
those completed results is counted here as a new theorem.

The present work concentrates on places where the existing hypotheses
leave a real mathematical boundary: outer orders below one, simultaneous
growth of asymptotic indices, and the difference between a set of
linear symmetry rows and its multiplicative closure. These boundaries
lead to the principal results in Table~\ref{tab:results}.

\begin{table}[htbp]
\centering\small
\begin{tabular}{@{}p{0.25\linewidth}p{0.67\linewidth}@{}}
\toprule
Result & Precise advance\\
\midrule
Sharp finite-measure domain & A necessary and sufficient condition for
$H_{n-1}^{(b)}/n^a$ to have a finite signed moment measure; explicit
interior densities and critical atoms
(Theorem~\ref{subunit:thm:domain}).\\[5pt]
Complex and angular zeros & Nonvanishing on the principal slit plane and
one simple imaginary-part zero on each open upper arc throughout the
full measure domain (Theorem~\ref{subunit:thm:zerofree}).\\[5pt]
Endpoint scales & Two power laws, an exponential logarithmic corner,
weak limits, and a sharp divergent signed-mass scale
(Theorem~\ref{subunit:thm:boundarylayer} and
Proposition~\ref{subunit:prop:otheredges}).\\[5pt]
Existing radial conjecture & The exact sign of its quadratic radial
coefficient for every real $a\ge2,b>0$, its $a=1$ boundary, and a
unique fractional threshold (Theorem~\ref{thm:radius-local-sign});
proved local maxima and minima at two certified inner orders
(Theorems~\ref{turn:thm:maximum} and \ref{turn:thm:minimum}).\\[5pt]
Adjacent-index zeros & A complete Bessel expansion for $n=k+r$, every
fixed integer $r$, uniform in $0\le\rho\le1$, with explicit zero
coefficients and exponential deformation bounds
(Theorems~\ref{bessel:thm:profile}--\ref{bessel:thm:zeros}).\\[5pt]
Small Lerch parameter & Complete positive-zero counts and central
Puiseux branches for every fixed $(n,k)$
(Theorem~\ref{unfold:thm:count}).\\[5pt]
Growing reflected moments & A uniform transition at
$n=(m+1)(\log m+c)$, its first correction and lattice location
(Theorem~\ref{thm:gamma-transition}).\\[5pt]
Cayley quotient and identities & Exact all-weight raw ranks and
shuffle-ideal Hilbert characters; a proved explicit family containing
$S_6$ (Theorems~\ref{cayley:thm:raw}--\ref{cayley:thm:ideal} and
Proposition~\ref{cayley:prop:harmonic}).\\
\bottomrule
\end{tabular}
\caption{Theorem-level advances relative to the inspected source.
Finite computations audit these results; their all-parameter statements
are established by the written proofs.}
\label{tab:results}
\end{table}

\subsection{Relation to established mathematics}
The regularized Hurwitz and digamma integrals used below are classical
\cite{DLMFhurwitz,DLMFdigamma}. The Bessel limit is proved here directly
from a finite elementary-symmetric product; the Bessel series and simple
positive zeros are classical \cite{DLMFbessel}. Its form is reminiscent
of Mehler--Heine limits, but no identification of these coefficient
polynomials with a classical orthogonal family is assumed.
The shuffle polynomial theorem is due to the classical work of Radford
and Melan\c{c}on--Reutenauer
\cite{cayley:radford,cayley:melancon-reutenauer}.
The Cayley transformation itself is inherited from the manuscript's
convergent octahedral argument; it is not introduced here as a new
change of variables. The broader use of octahedral symmetry at level
four has a precedent in Zhao \cite{ZhaoOctahedral}.

Reflected products of powers of $\log\Gamma$ were studied by
Amdeberhan, Coffey, Espinosa, Koutschan, Manna and Moll
\cite{ACEKMM}. In particular, the moments, their reflection symmetry,
and elementary identities obtained from Euler's reflection formula are
not new. Our contribution is the joint growth regime and its controlled
correction. Targeted primary-source research supports these attributions;
it does not establish exhaustive historical priority for every new
deduction in this article.

\section{Conventions and inherited analytic framework}
\label{sec:conventions}

For real $b>0$, let
\[
 H_N^{(b)}=\sum_{j=1}^N j^{-b},\qquad H_0^{(b)}=0,
 \qquad H_N=H_N^{(1)}.
\]
We use decreasing series indices and the outer-first convention
\[
 \Li_{s_1,\ldots,s_d}(z_1,\ldots,z_d)
 =\sum_{n_1>\cdots>n_d\ge1}
 \frac{z_1^{n_1}\cdots z_d^{n_d}}
 {n_1^{s_1}\cdots n_d^{s_d}}
\]
where convergent, followed by the stated analytic continuation.
In particular,
\[
 F_{a,b}(z)=\Li_{a,b}(z,1)
 =\sum_{n\ge1}\frac{H_{n-1}^{(b)}}{n^a}z^n,
 \qquad |z|<1.
\]
The slit plane is $\mathbb C\setminus[1,\infty)$, and logarithms
there are continued from the disk. Later, $g_{a,b}=\Im F_{a,b}(i)$
denotes a Gaussian double, $\beta(s)=\sum_{j\ge0}(-1)^j/(2j+1)^s$
is the Dirichlet beta function, and
\[
 S_p=\sum_{j\ge0}\frac{(-1)^jH_j}{(2j+1)^p}.
\]
The last sum is absolute for integers $p>1$ and converges conditionally
for $p=1$. Throughout, $\gamma$ without an index is Euler's constant.

The Stieltjes constants use
\[
 \zeta(s,a)=\frac1{s-1}
  +\sum_{n\ge0}\frac{(-1)^n\gamma_n(a)}{n!}(s-1)^n.
\]
A superscript $(k)$ on $\gamma_n$ means differentiation in $a$.
For $0\le\rho<1$ set
\[
 \ell_n(\rho,a)=\sum_{j\ge0}\frac{\rho^j\log^n(a+j)}{a+j},
 \qquad
 \mathcal D_{n,k}^{\rho}(a)=\partial_a^k\ell_n(\rho,a),
 \qquad k\ge1,
\]
and define $\mathcal D_{n,k}^{1}(a)=\gamma_n^{(k)}(a)$.
Only the differentiated family is unified at $\rho=1$; the series
defining $\ell_n$ generally diverges there. The sign-normalized family is
\[
 F_{n,k}^{\rho}(a)=(-1)^{n+k}\mathcal D_{n,k}^{\rho}(a).
\]
The same symbol $F$ with different indices is conventional in the
source: $F_{a,b}(z)$ is a double polylogarithm, while
$F_{n,k}^{\rho}(a)$ is the Stieltjes--Lerch derivative family.
The radius $\rho$ in the angular sections and the Lerch parameter
$\rho$ in the zero sections are separate variables.

The master coefficient formula, reproduced in
\eqref{bessel:eq:shift}, follows from differentiating
$(a+j)^{-s}$ in $a$ and extracting the coefficient of $s-1$.
For $k\ge1$ its shift series and every fixed number of $a$
derivatives converge locally uniformly for $a>0$, uniformly in
$0\le\rho\le1$. The bound is a fixed power of
$\log(a+j)$ divided by $(a+j)^{k+1}$.
This convergence justifies the classical endpoint without assigning
a value to the divergent undifferentiated series.

All finite signed measures in the article are real measures of finite
total variation. Weak convergence means convergence against every
continuous function on $[0,1]$. Big-$O$ constants may depend on fixed
indices and fixed compact parameter sets; every claimed uniformity is
stated explicitly.
